
% -------------------------------------------------
% فصل 3
% -------------------------------------------------

\projectchapter{داده‌ها و پیش‌پردازش}{
در این فصل، دو مجموعه‌داده مستقل متنی و گفتاری پروژه معرفی شده و مراحل عملی آماده‌سازی آن‌ها تشریح می‌شود. ابتدا ساختار و توزیع داده‌های متنی، نحوه پاک‌سازی و تبدیل آن‌ها به فضای ویژگی بررسی شده و سپس فرایندی که برای استانداردسازی فایل‌های صوتی، تولید فراداده و استخراج ویژگی‌های گفتاری انجام شده است ارائه می‌شود. هدف این فصل، تشریح روند تبدیل داده‌های خام به ورودی‌های قابل استفاده برای مدل‌های یادگیری ماشین است.
}


% =================================================
\section{مقدمه}
% =================================================

کیفیت داده ورودی یکی از عوامل مهم در عملکرد مدل‌های یادگیری ماشین است. حتی در صورت انتخاب یک الگوریتم مناسب، وجود داده‌های ناسازگار، ناقص یا دارای ساختار نامناسب می‌تواند بر نتیجه نهایی تأثیر منفی بگذارد. به همین دلیل، پیش از مرحله آموزش مدل لازم است داده‌های خام بررسی، پاک‌سازی و به یک نمایش مناسب تبدیل شوند.

در پروژه حاضر، دو مجموعه‌داده با ساختارهای متفاوت مورد استفاده قرار گرفته است. بخش نخست شامل فایل‌های متنی در چند زبان مختلف و بخش دوم شامل داده‌های گفتاری است که پس از پاک‌سازی صوت و استخراج ویژگی در قالب جدول عددی ذخیره شده‌اند.

از آنجا که ماهیت این دو نوع داده متفاوت است، مراحل پیش‌پردازش آن‌ها نیز به صورت مستقل طراحی شده است. با این حال، در هر دو بخش هدف نهایی یکسان است: تبدیل داده خام به ساختاری عددی، منظم و قابل استفاده در الگوریتم‌های یادگیری ماشین.

مبانی نظری روش‌های استخراج ویژگی در فصل دوم تشریح شد؛ بنابراین در این فصل تمرکز اصلی بر داده‌های واقعی، نحوه پیاده‌سازی و پارامترهایی است که در نسخه نهایی پروژه مورد استفاده قرار گرفته‌اند. مراحل پاک‌سازی و استخراج ویژگی صوتی پیش‌تر توسط نگارنده اجرا و فایل‌های
\lr{features.csv}
و
\lr{metadata.csv}
تولید شده‌اند؛ خط لوله مدل‌سازی فعلی این دو فایل را به‌عنوان ورودی فقط‌خواندنی دریافت می‌کند و استخراج ویژگی صوتی را دوباره اجرا نمی‌کند.


% =================================================
\section{ساختار کلی داده‌های پروژه}
% =================================================

داده‌های پروژه در دو مجموعه مستقل سازمان‌دهی شده‌اند:

\begin{itemize}

	\item یک مجموعه مستقل از فایل‌های متنی برای شناسایی زبان از متن؛

	\item یک مجموعه گفتاری مستقل که فایل‌های صوتی پاک‌سازی‌شده، فراداده و جدول ویژگی‌های استخراج‌شده را در بر می‌گیرد.

\end{itemize}

در بخش متن، هر نمونه شامل یک فایل متنی، برچسب زبان و فراداده اختیاری استخراج‌شده از ساختار مسیر است. پس از بارگذاری و پاک‌سازی، متن‌ها به بردارهای عددی تبدیل می‌شوند.

در بخش گفتار نیز هر سطر از جدول نهایی با یک فایل صوتی متناظر است و شامل مسیر فایل، برچسب زبان، اطلاعات جنسیت گوینده و ویژگی‌های عددی استخراج‌شده است.

دو مجموعه داده به صورت مستقل پردازش شده‌اند. هیچ تطبیق سطربه‌سطر میان نمونه‌های متن و گفتار فرض نشده و هیچ ترکیب مستقیمی میان ویژگی‌های دو بخش در مرحله مدل‌سازی انجام نشده است؛ بنابراین نتایج آن‌ها نیز به‌عنوان یک آزمایش زوج‌شده تفسیر نمی‌شوند.


% =================================================
\section{آماده‌سازی داده‌های متنی}
% =================================================


% -------------------------------------------------
\subsection{ساختار و توزیع داده‌های متنی}
% -------------------------------------------------

مجموعه داده متنی نهایی شامل
\textbf{608 نمونه}
از شش زبان آلمانی، اسپانیایی، ایتالیایی، کره‌ای، پرتغالی و ژاپنی است.

ساختار پوشه‌های داده به شکلی طراحی شده است که هر زبان دارای یک کد اختصاصی باشد. نگاشت مورد استفاده در پروژه در جدول زیر ارائه شده است.

\begin{table}[H]
\centering
\caption{توزیع نمونه‌های مجموعه داده متنی}
\label{tab:text-data-distribution}
\begin{tabular}{c c c c}
\toprule
\textbf{زبان} &
\textbf{کد پوشه} &
\textbf{تعداد نمونه} &
\textbf{میانگین طول متن} \\
\midrule
آلمانی     & \lr{de} & \lr{104} & \lr{313.5} \\
اسپانیایی  & \lr{es} & \lr{100} & \lr{313.8} \\
ایتالیایی  & \lr{it} & \lr{106} & \lr{313.0} \\
کره‌ای     & \lr{ko} & \lr{98}  & \lr{318.3} \\
پرتغالی    & \lr{po} & \lr{100} & \lr{314.9} \\
ژاپنی      & \lr{ja} & \lr{100} & \lr{315.0} \\
\midrule
\textbf{مجموع} & \lr{--} & \textbf{\lr{608}} & \lr{--} \\
\bottomrule
\end{tabular}
\end{table}

همان‌طور که جدول \ref{tab:text-data-distribution} نشان می‌دهد، تعداد نمونه‌ها بین زبان‌های مختلف تقریباً متوازن است و اختلاف تعداد نمونه‌های کلاس‌ها محدود است.

میانگین طول متن در این جدول پس از یکسان‌سازی فاصله‌ها و برحسب تعداد کاراکتر محاسبه شده است.

این ویژگی اهمیت دارد، زیرا عدم توازن شدید میان کلاس‌ها می‌تواند باعث شود مدل در فرایند آموزش به سمت زبان‌هایی که تعداد نمونه بیشتری دارند متمایل شود.


% -------------------------------------------------
\subsection{بارگذاری و برچسب‌گذاری داده‌های متنی}
% -------------------------------------------------

برای بارگذاری داده‌ها از مسیر
\lr{\texttt{data/raw/text}}
استفاده شده است. تمام فایل‌های دارای پسوند
\lr{\texttt{.txt}}
به صورت بازگشتی از شش پوشه زبانی تعریف‌شده در فایل تنظیمات خوانده شده‌اند. نام پوشه نخست، کد زبان را مشخص می‌کند و سپس این کد به نام کامل زبان نگاشت می‌شود.

هر فایل متنی با کدگذاری
\lr{UTF-8}
و به‌صورت سخت‌گیرانه خوانده شده است؛ بنابراین در صورت رخ‌دادن خطای رمزگشایی، اجرا با پیام خطا متوقف می‌شود و کاراکتر نامعتبر به‌صورت خاموش حذف نمی‌شود. در اجرای نهایی، هر 608 فایل با موفقیت رمزگشایی شدند و هیچ کاراکتر جایگزین یونیکد در متن پاک‌سازی‌شده مشاهده نشد.

برای هر نمونه اطلاعات زیر ذخیره شده است:

\begin{itemize}

	\item شناسه نمونه شامل مسیر نسبی فایل؛

	\item زبان نمونه؛

	\item اطلاعات جنسیت استخراج‌شده از پوشه دوم، در صورت وجود؛

	\item متن خام و متن پاک‌سازی‌شده؛

	\item شناسه گروه مربوط به گوینده یا منبع نمونه.

\end{itemize}


% -------------------------------------------------
\subsection{پاک‌سازی متون}
% -------------------------------------------------

هدف اصلی مرحله پاک‌سازی، ایجاد یک نمایش یکنواخت از متن بدون حذف اطلاعاتی است که ممکن است برای تشخیص زبان مفید باشند.

پاک‌سازی مورد استفاده در این پروژه به صورت کنترل‌شده و حداقلی انجام شده است.

در این مرحله، تمام فاصله‌ها، خطوط جدید و
\lr{Tab}
های متوالی به یک فاصله واحد تبدیل شده و فاصله‌های ابتدا و انتهای متن حذف شده‌اند.

به صورت مفهومی عملیات انجام‌شده را می‌توان به شکل زیر نمایش داد:

\begin{center}
\lr{\texttt{Multiple Whitespaces}}
$\longrightarrow$
\lr{\texttt{Single Space}}
\end{center}

برخلاف برخی مسائل پردازش متن، در این پروژه کاراکترهای خاص زبان‌ها حذف نشده‌اند. همچنین علائم نگارشی، اعداد، حروف دارای
\lr{Accent}،
\lr{Umlaut}
و کاراکترهای زبان‌های کره‌ای و ژاپنی حفظ شده‌اند.

در این مرحله، تبدیل حروف به حالت کوچک، ریشه‌یابی، بن‌واژه‌سازی، نویسه‌گردانی یا توکن‌سازی ویژه زبان انگلیسی انجام نشده است. تبدیل حروف لاتین به حالت کوچک بعداً درون بردارسازهای
\lr{TF-IDF}
اعمال می‌شود.

دلیل این تصمیم آن است که برخی از این کاراکترها می‌توانند اطلاعات ارزشمندی درباره زبان یک نمونه در اختیار سیستم قرار دهند.

پس از انجام پاک‌سازی، طول متن‌ها نیز برحسب تعداد کاراکتر محاسبه شده است.

در مرحله کنترل کیفیت مشخص شد که هیچ نمونه‌ای پس از پاک‌سازی خالی نشده است؛ بنابراین:

\begin{center}
تعداد نمونه قبل از پاک‌سازی = 608
\end{center}

و

\begin{center}
تعداد نمونه بعد از پاک‌سازی = 608
\end{center}

یکسان‌سازی فاصله‌ها محتوای 14 نمونه را تغییر داد. سیاست پروژه برای متن‌های کاملاً تکراری این است که پس از یکسان‌سازی فاصله‌ها، اولین مسیر براساس ترتیب واژگانی نگهداری و موارد بعدی حذف شوند؛ بااین‌حال، در اجرای نهایی هیچ متن کاملاً تکراری حذف نشد.


% -------------------------------------------------
\subsection{تبدیل متون به فضای ویژگی}
% -------------------------------------------------

پس از پاک‌سازی، متون به بردارهای عددی تبدیل شده‌اند.

همان‌طور که در فصل دوم تشریح شد، برای استخراج اطلاعات از دو سطح متفاوت متن، دو بردارساز مستقل در نظر گرفته شده است.

برای بخش کاراکتری تنظیمات زیر مورد استفاده قرار گرفته‌اند:

\begin{itemize}

	\item نوع تحلیل:
	\lr{\texttt{char\_wb}}؛

	\item بازه
	\lr{N-Gram}:
	2 تا 4؛

	\item حداکثر تعداد ویژگی:
	3000؛

	\item
	\lr{\texttt{sublinear\_tf=True}}؛

	\item
	\lr{\texttt{min\_df=1}}؛

	\item تبدیل حروف لاتین به حالت کوچک با
	\lr{\texttt{lowercase=True}}.

\end{itemize}

برای ویژگی‌های سطح کلمه نیز تنظیمات زیر استفاده شده‌اند:

\begin{itemize}

	\item نوع تحلیل:
	\lr{\texttt{word}}؛

	\item بازه
	\lr{N-Gram}:
	1 تا 2؛

	\item حداکثر تعداد ویژگی:
	2000؛

	\item
	\lr{\texttt{sublinear\_tf=True}}؛

	\item
	\lr{\texttt{min\_df=1}}؛

	\item
	\lr{\texttt{lowercase=True}}.

\end{itemize}

دو مجموعه ویژگی با استفاده از
\lr{FeatureUnion}
با یکدیگر ترکیب شده‌اند. در مسیر طبقه‌بندی، بردارسازها فقط روی بخش آموزشی برازش شده‌اند. در اجرای نهایی، واژگان حاصل شامل 5000 ویژگی بوده است که از سقف 3000 ویژگی کاراکتری و 2000 ویژگی واژگانی تشکیل می‌شود.

ابعاد ماتریس‌های واقعی تولیدشده برای مسیر طبقه‌بندی به صورت زیر بوده است:

\begin{table}[H]
\centering
\caption{ابعاد فضای ویژگی متنی}
\label{tab:text-feature-shapes}
\begin{tabular}{c c}
\toprule
\textbf{فضای ویژگی} & \textbf{ابعاد ماتریس} \\
\midrule
ترکیب ویژگی‌های آموزشی پیش از انتخاب ویژگی & $397 \times 5000$ \\
ترکیب ویژگی‌های آزمون پیش از انتخاب ویژگی & $211 \times 5000$ \\
ویژگی‌های آموزشی پس از انتخاب ویژگی & $397 \times 50$ \\
ویژگی‌های آزمون پس از انتخاب ویژگی & $211 \times 50$ \\
\bottomrule
\end{tabular}
\end{table}

برای مسیر خوشه‌بندی، انتخاب ویژگی نظارت‌شده
\lr{Chi-Square}
استفاده نشده است. بردارسازی روی هر 608 متن بدون استفاده از برچسب زبان انجام شده و سپس با
\lr{Truncated SVD}
به 50 مؤلفه کاهش یافته است. در پایان نیز هر بردار با نرم
\lr{L2}
نرمال‌سازی شده و ماتریسی با ابعاد
$608\times50$
برای خوشه‌بندی ایجاد شده است.


% -------------------------------------------------
\subsection{آماده‌سازی برچسب‌ها و انتخاب ویژگی}
% -------------------------------------------------

برچسب زبان‌ها در مسیر طبقه‌بندی متن به صورت رشته‌ای نگهداری شده‌اند و مدل‌های
\lr{scikit-learn}
مستقیماً با همین برچسب‌ها آموزش می‌بینند؛ بنابراین در نسخه نهایی بخش متن، مرحله جداگانه‌ای برای
\lr{LabelEncoder}
وجود ندارد.

کلاس‌های نهایی عبارت‌اند از:

\begin{center}
\lr{
german,\;
italian,\;
japanese,\;
korean,\;
portuguese,\;
spanish
}
\end{center}

در مرحله بعد، فضای ترکیبی شامل 5000 ویژگی با استفاده از
\lr{SelectKBest}
و معیار
\lr{Chi-Square}
کاهش یافته است.

تعداد ویژگی‌های انتخاب‌شده برابر با:

\begin{equation}
K_{\mathrm{best}}=50
\end{equation}

این انتخاب‌گر فقط روی بخش آموزشی برازش شده است. در نتیجه، ماتریس‌های نهایی آموزش و آزمون دارای ابعاد زیر هستند:

\begin{equation}
X_{\mathrm{text,train}}
\in
\mathbb{R}^{397\times50},
\qquad
X_{\mathrm{text,test}}
\in
\mathbb{R}^{211\times50}
\end{equation}

این کاهش ابعاد باعث می‌شود تنها ویژگی‌هایی که در داده آموزشی ارتباط بیشتری با برچسب زبان دارند برای طبقه‌بندی نگهداری شوند. به دلیل وابستگی این روش به برچسب، 50 ویژگی منتخب در خوشه‌بندی متن استفاده نشده‌اند.


% -------------------------------------------------
\subsection{جلوگیری از نشت اطلاعات در بخش متن}
% -------------------------------------------------

یکی از نکات مهم در فرایند ارزیابی مدل، جلوگیری از
\lr{Data Leakage}
است.

اگر بردارساز
\lr{TF-IDF}
یا روش انتخاب ویژگی پیش از تقسیم داده‌های آموزش و آزمون روی کل مجموعه داده برازش شود، اطلاعات مربوط به داده آزمون می‌تواند به صورت غیرمستقیم وارد فرایند آموزش شود و نتیجه ارزیابی بیش از حد خوش‌بینانه باشد.

برای تشکیل گروه‌ها، شناسه 9 رقمی ابتدای نام فایل استخراج و با نام زبان ترکیب شده است. فایل‌هایی که نام آن‌ها با این الگو منطبق نیست، هرکدام یک گروه یکتا دریافت کرده‌اند. در داده نهایی 284 گروه وجود داشت که 44 گروه بیش از یک نمونه داشتند.

سپس با استفاده از نخستین بخش یک تقسیم‌بندی سه‌قسمتی
\lr{Stratified Group K-Fold}،
397 نمونه به مجموعه آموزش و 211 نمونه به مجموعه آزمون اختصاص یافت. همه زبان‌ها در هر دو بخش حضور دارند و تعداد گروه‌های مشترک میان آموزش و آزمون صفر است.

در بخش طبقه‌بندی، مراحل پاک‌سازی، استخراج
\lr{TF-IDF}
و انتخاب ویژگی
\lr{Chi-Square}
درون خط لوله مدل قرار گرفته‌اند. این اجزا در هر بار آموزش و هر دور اعتبارسنجی فقط روی داده آموزشی
\lr{fit}
می‌شوند و برای داده اعتبارسنجی یا آزمون فقط
\lr{transform}
انجام می‌شود.

همچنین پیش از آموزش مدل، وجود متن‌های یکسان پس از نرمال‌سازی میان مجموعه آموزش و آزمون بررسی شد. تعداد متن‌های مشترک و نیز تعداد گروه‌های مشترک در اجرای نهایی هر دو برابر صفر بود.


% -------------------------------------------------
\subsection{خروجی‌های مرحله آماده‌سازی متن}
% -------------------------------------------------

برای امکان استفاده مجدد از نتایج، خروجی‌های اصلی مرحله پیش‌پردازش و استخراج ویژگی ذخیره شده‌اند.

مهم‌ترین آن‌ها عبارت‌اند از:

\begin{itemize}

	\item \lr{\texttt{artifacts/text/audit/text\_audit.json}} برای خلاصه کشف داده، کنترل کیفیت، توزیع کلاس‌ها و اطلاعات تقسیم داده؛

	\item \lr{\texttt{artifacts/text/audit/class\_distribution.csv}} برای توزیع زبان‌ها؛

	\item \lr{\texttt{artifacts/text/preprocessing/cleaning\_counts.csv}} برای تعداد نمونه‌ها در مراحل پاک‌سازی؛

	\item \lr{\texttt{artifacts/text/splits/holdout\_split.csv}} برای ثبت گروه و بخش هر نمونه؛

	\item \lr{\texttt{artifacts/text/features/training\_vectorizer.joblib}} برای خط لوله پاک‌سازی و بردارسازی برازش‌شده روی داده آموزشی؛

	\item \lr{\texttt{artifacts/text/features/training\_feature\_selector.joblib}} برای انتخاب‌گر 50 ویژگی برازش‌شده روی داده آموزشی؛

	\item \lr{\texttt{artifacts/text/features/feature\_statistics.json}} برای تنظیمات و ابعاد واقعی ماتریس‌ها.

\end{itemize}

در نسخه نهایی، ماتریس‌های
\lr{TF-IDF}
کل داده به صورت فایل‌های مستقل
\lr{NPZ}
ذخیره نشده‌اند؛ زیرا برای طبقه‌بندی باید درون هر خط لوله و فقط از داده آموزشی ساخته شوند. نمایش بدون برچسب مخصوص خوشه‌بندی نیز همراه با مدل کاهش ابعاد در بخش مصنوعات خوشه‌بندی ذخیره شده است.


% =================================================
\section{آماده‌سازی داده‌های گفتاری}
% =================================================


% -------------------------------------------------
\subsection{ساختار و توزیع داده‌های صوتی}
% -------------------------------------------------

داده‌های این بخش حاصل فرایند پاک‌سازی و استخراج ویژگی‌ای هستند که نگارنده روی فایل‌های صوتی اجرا کرده است. اسکریپت‌ها، دفترچه استخراج ویژگی و فایل‌های خروجی در مسیر
\lr{\texttt{Speech\_Identification/clean data and data extraction}}
نگهداری می‌شوند. در اجرای آزمایش‌های نسخه نهایی، این مسیر تغییر داده نشده و جدول ویژگی موجود مستقیماً بارگذاری شده است.

مجموعه داده گفتاری نهایی شامل
\textbf{720 فایل صوتی}
از چهار زبان آلمانی، اسپانیایی، ایتالیایی و کره‌ای است.

تعداد نمونه‌های هر زبان کاملاً برابر بوده و برای هر زبان 180 فایل در نظر گرفته شده است.

\begin{table}[H]
\centering
\caption{توزیع نمونه‌های صوتی براساس زبان}
\label{tab:speech-language-distribution}
\begin{tabular}{c c c}
\toprule
\textbf{زبان} & \textbf{برچسب} & \textbf{تعداد نمونه} \\
\midrule
آلمانی    & \lr{de} & \lr{180} \\
اسپانیایی & \lr{es} & \lr{180} \\
ایتالیایی & \lr{it} & \lr{180} \\
کره‌ای    & \lr{ko} & \lr{180} \\
\midrule
\textbf{مجموع} & \lr{--} & \textbf{\lr{720}} \\
\bottomrule
\end{tabular}
\end{table}

داده‌ها از نظر جنسیت گوینده نیز متوازن هستند:

\begin{table}[H]
\centering
\caption{توزیع نمونه‌های صوتی براساس جنسیت گوینده}
\label{tab:speech-gender-distribution}
\begin{tabular}{c c}
\toprule
\textbf{جنسیت} & \textbf{تعداد نمونه} \\
\midrule
\lr{Female} & \lr{360} \\
\lr{Male}   & \lr{360} \\
\midrule
\textbf{مجموع} & \textbf{\lr{720}} \\
\bottomrule
\end{tabular}
\end{table}

توازن کامل کلاس‌های زبان باعث می‌شود مقایسه عملکرد مدل روی زبان‌های مختلف تحت تأثیر اختلاف شدید تعداد نمونه‌ها قرار نگیرد.


% -------------------------------------------------
\subsection{استانداردسازی فایل‌های صوتی}
% -------------------------------------------------

فایل‌های صوتی خام در زمان آماده‌سازی داده گفتار، در مسیر
\lr{\texttt{data/raw}}
نسبت به پوشه اسکریپت پاک‌سازی قرار داشته و عمدتاً دارای قالب
\lr{MP3}
بوده‌اند.

برای ایجاد یک ساختار صوتی یکسان، اسکریپت
\lr{PowerShell}
تمام فایل‌های صوتی را به صورت بازگشتی پیمایش کرده و با استفاده از
\lr{FFmpeg}
نسخه استاندارد آن‌ها را در مسیر
\lr{\texttt{data/clean}}
تولید کرده است.

ساختار نسبی پوشه‌ها نیز در این فرایند حفظ شده است.

تنظیمات فعال مورد استفاده در تبدیل فایل‌های صوتی شامل موارد زیر است:

\begin{itemize}

	\item تبدیل فایل‌ها از
	\lr{MP3}
	به
	\lr{WAV}؛

	\item تبدیل صدا به حالت تک‌کاناله با
	\lr{\texttt{"-ac 1"}}؛

	\item بازنمونه‌برداری با نرخ
	\lr{16 kHz}؛

	\item استفاده از کدگذاری
	\lr{\texttt{"pcm\_s16le"}}؛

	\item نرمال‌سازی شدت صوت با فیلتر
	\lr{loudnorm}.

\end{itemize}

پارامترهای نرمال‌سازی شدت صوت در نسخه نهایی به صورت زیر تنظیم شده‌اند:

\begin{center}
\lr{\texttt{I=-16,\ TP=-1.5,\ LRA=11}}
\end{center}

که هدف از آن کاهش اختلاف شدید سطح صوت میان فایل‌های مختلف و ایجاد شرایط یکنواخت‌تر برای استخراج ویژگی است.

در نسخه نهایی زنجیره پردازش صوت، عملیات حذف سکوت از ابتدا و انتهای فایل به فرمان فعال
\lr{FFmpeg}
اضافه نشده است؛ بنابراین مراحل اصلی پیش‌پردازش شامل تبدیل قالب، تک‌کاناله‌سازی، بازنمونه‌برداری و نرمال‌سازی شدت صوت هستند.

این عملیات مرحله ایجاد داده بوده است و در خط لوله جاری آموزش مدل دوباره اجرا نمی‌شود؛ ورودی بخش گفتار از فایل ویژگی ازپیش‌استخراج‌شده آغاز می‌شود.


% -------------------------------------------------
\subsection{ایجاد فایل فراداده}
% -------------------------------------------------

پس از آماده‌سازی فایل‌های صوتی، برای مدیریت سیستماتیک نمونه‌ها یک فایل
\lr{Metadata}
ایجاد شده است.

ساختار مسیر فایل‌ها به صورت کلی مشابه الگوی زیر است:

\begin{center}
\lr{\texttt{data/clean/<language>/<gender>/<file>.wav}}
\end{center}

اسکریپت تولید فراداده، زبان و جنسیت را مستقیماً از ساختار مسیر استخراج کرده و برای هر نمونه سه ستون زیر ایجاد می‌کند:

\begin{itemize}

	\item \lr{\texttt{filepath}}: مسیر نسبی فایل صوتی؛

	\item \lr{\texttt{lang}}: برچسب زبان؛

	\item \lr{\texttt{gender}}: جنسیت گوینده.

\end{itemize}

فایل نهایی
\lr{\texttt{metadata.csv}}
دارای 720 سطر و 3 ستون است و هیچ مقدار گمشده‌ای در آن مشاهده نشده است.

استفاده از مسیر نسبی به جای مسیر کامل سیستم نیز باعث می‌شود ساختار پروژه قابلیت جابه‌جایی بیشتری میان سیستم‌های مختلف داشته باشد.

هنگام بارگذاری داده در خط لوله فعلی، فراداده براساس ستون
\lr{\texttt{filepath}}
با
\lr{features.csv}
هم‌تراز می‌شود و برابری مقادیر مشترک کنترل می‌گردد؛ بنابراین اتصال دو فایل به ترتیب تصادفی سطرها وابسته نیست.


% -------------------------------------------------
\subsection{پیاده‌سازی استخراج ویژگی‌های صوتی}
% -------------------------------------------------

پس از تولید فراداده، در دفترچه استخراج ویژگی هر فایل صوتی با استفاده از کتابخانه
\lr{Librosa}
خوانده شده است.

تنظیمات اصلی مرحله بارگذاری عبارت‌اند از:

\begin{itemize}

	\item نرخ نمونه‌برداری:
	\lr{16000 Hz}؛

	\item تبدیل به سیگنال تک‌کاناله:
	\lr{\texttt{mono=True}}؛

	\item تعداد ضرایب پایه
	\lr{MFCC}:
	\lr{20}.

\end{itemize}

برای تبدیل ویژگی‌های وابسته به زمان به یک بردار با طول ثابت، میانگین و انحراف معیار آن‌ها در طول فریم‌های فایل صوتی محاسبه شده است.

ویژگی‌های نهایی از گروه‌های زیر تشکیل شده‌اند:

\begin{table}[H]
\centering
\caption{ساختار ویژگی‌های استخراج‌شده از هر فایل صوتی}
\label{tab:speech-feature-count}
\begin{tabular}{c c}
\toprule
\textbf{گروه ویژگی} & \textbf{تعداد ویژگی نهایی} \\
\midrule
\lr{MFCC} & \lr{40} \\
\lr{Delta MFCC} & \lr{40} \\
\lr{Delta-Delta MFCC} & \lr{40} \\
\lr{Zero Crossing Rate} & \lr{2} \\
\lr{Spectral Centroid} & \lr{2} \\
\lr{Spectral Rolloff} & \lr{2} \\
\lr{Spectral Bandwidth} &\lr{2} \\
\lr{Spectral Flatness} & \lr{2} \\
\lr{RMS Energy} & \lr{2} \\
\lr{Duration} &\lr{1} \\
\midrule
\textbf{مجموع} & \textbf{\lr{133}} \\
\bottomrule
\end{tabular}
\end{table}

بنابراین برای هر فایل صوتی یک بردار شامل 133 مقدار عددی ایجاد شده است.

پس از اضافه شدن سه ستون توصیفی
\lr{\texttt{filepath}},
\lr{\texttt{lang}}
و
\lr{\texttt{gender}},
فایل نهایی دارای 136 ستون است.

این فایل با نام
\lr{\texttt{features.csv}}
ذخیره شده و نقطه شروع بخش گفتار در خط لوله آزمایش‌های فعلی است.


% -------------------------------------------------
\subsection{کنترل کیفیت ویژگی‌های صوتی}
% -------------------------------------------------

استخراج ویژگی روی تمام 720 نمونه صوتی انجام شده است.

کد استخراج ویژگی به شکلی طراحی شده است که در صورت بروز خطا برای یک فایل، مسیر آن فایل و متن خطا ذخیره شده و نمونه معیوب از مجموعه نهایی حذف شود.

در اجرای نهایی پروژه:

\begin{itemize}

	\item تعداد فایل‌های ورودی برابر 720 بود؛

	\item تعداد فایل‌های پردازش‌شده موفق برابر 720 بود؛

	\item هیچ نمونه‌ای به علت خطای استخراج ویژگی حذف نشد؛

	\item تعداد مقادیر
	\lr{NaN}
	در فایل نهایی برابر صفر بود.

	\item تعداد مقادیر بی‌نهایت، سطرهای کاملاً تکراری و شناسه‌های فایل تکراری برابر صفر بود؛

	\item هیچ ویژگی ثابت یا ویژگی نزدیک به ثابت با آستانه 99 درصد مشاهده نشد.

\end{itemize}

در نتیجه ابعاد فایل ویژگی نهایی برابر است با:

\begin{equation}
720\times136
\end{equation}

که 133 ستون آن ویژگی عددی هستند.

برای کنترل وابستگی نمونه‌های یک گوینده یا منبع، شناسه گروه از 9 رقم ابتدای نام فایل استخراج شد. هر 720 مسیر با این الگو منطبق بودند و در مجموع 78 گروه به دست آمد. تقسیم نگه‌داشته‌شده با نخستین بخش یک
\lr{Stratified Group K-Fold}
پنج‌قسمتی انجام شد و شامل 572 نمونه آموزشی و 148 نمونه آزمون بود؛ هیچ گروه مشترکی میان دو بخش وجود نداشت.


% -------------------------------------------------
\subsection{بررسی مدت زمان نمونه‌های صوتی}
% -------------------------------------------------

از آنجا که طول فایل‌های صوتی یکسان نیست، مدت زمان هر فایل نیز محاسبه و به عنوان یکی از ویژگی‌های مجموعه داده ذخیره شده است.

آمار کلی مدت زمان فایل‌ها در جدول زیر ارائه شده است.

\begin{table}[H]
\centering
\caption{آمار مدت زمان نمونه‌های صوتی}
\label{tab:audio-duration-statistics}
\begin{tabular}{c c}
\toprule
\textbf{شاخص} & \textbf{مدت زمان برحسب ثانیه} \\
\midrule
تعداد نمونه & \lr{720} \\
میانگین & \lr{63.19} \\
انحراف معیار & \lr{55.22} \\
کمینه & \lr{39.00} \\
میانه & \lr{60.68} \\
چارک اول & \lr{58.00} \\
چارک سوم & \lr{63.91} \\
بیشینه & \lr{1534.25} \\
\bottomrule
\end{tabular}
\end{table}

اکثر فایل‌های مجموعه داده دارای مدت زمانی در حدود یک دقیقه هستند. با این حال، یک نمونه از زبان ایتالیایی دارای مدت زمان تقریباً
1534
ثانیه است که نسبت به سایر نمونه‌ها بسیار طولانی‌تر محسوب می‌شود.

این نمونه در روند فعلی پروژه حذف نشده و در مجموعه نهایی باقی مانده است. از آنجا که مدت زمان نیز یکی از ویژگی‌های ورودی مدل محسوب می‌شود، وجود چنین مقدار دورافتاده‌ای باید هنگام تحلیل نتایج و اثر استانداردسازی ویژگی‌ها مورد توجه قرار گیرد.

ثبت و گزارش این مورد به عنوان بخشی از کنترل کیفیت داده، امکان تفسیر دقیق‌تر نتایج مدل را فراهم می‌کند.


% =================================================
\section{داده‌های نهایی آماده برای مدل‌سازی}
% =================================================

پس از تکمیل مراحل پیش‌پردازش، دو مجموعه داده نهایی برای بخش متن و گفتار ایجاد شده‌اند.

در بخش متنی، مجموعه داده پاک‌سازی‌شده شامل 608 نمونه از شش زبان است. برای ارزیابی طبقه‌بندی، 397 نمونه در بخش آموزش و 211 نمونه در بخش آزمون گروه‌ناهم‌پوشان قرار گرفته‌اند. فضای ویژگی
\lr{TF-IDF}
که فقط روی داده آموزشی برازش شده، 5000 بعد دارد و با انتخاب ویژگی
\lr{Chi-Square}
به 50 بعد کاهش یافته است.

بنابراین ورودی‌های مسیر طبقه‌بندی متن به صورت زیر قابل نمایش هستند:

\begin{equation}
X_{\mathrm{text,train}}
\in
\mathbb{R}^{397\times50},
\qquad
X_{\mathrm{text,test}}
\in
\mathbb{R}^{211\times50}
\end{equation}

در مسیر خوشه‌بندی متن، هر 608 نمونه بدون استفاده از برچسب زبان بردارسازی شده و پس از
\lr{Truncated SVD}
و نرمال‌سازی
\lr{L2}
نمایشی با ابعاد زیر ایجاد شده است:

\begin{equation}
X_{\mathrm{text,cluster}}
\in
\mathbb{R}^{608\times50}
\end{equation}

در بخش گفتار، فایل ویژگی دارای 720 نمونه و 136 ستون است.

سه ستون

\begin{center}
\lr{\texttt{filepath}},
\quad
\lr{\texttt{lang}},
\quad
\lr{\texttt{gender}}
\end{center}

به عنوان اطلاعات توصیفی در نظر گرفته می‌شوند.

برای تشکیل ورودی مدل، این سه ستون از ماتریس ویژگی حذف شده و 133 ستون عددی باقی می‌مانند:

\begin{equation}
X_{\mathrm{speech}}
\in
\mathbb{R}^{720\times133}
\end{equation}

برچسب هدف در بخش گفتار از ستون
\lr{\texttt{lang}}
استخراج می‌شود.

ستون
\lr{\texttt{gender}}
به عنوان ویژگی ورودی مدل استفاده نشده است و تنها برای توصیف و بررسی توازن مجموعه داده نگهداری می‌شود.

در ارزیابی طبقه‌بندی گفتار، 572 نمونه برای آموزش و 148 نمونه برای آزمون گروه‌ناهم‌پوشان استفاده شده‌اند. جایگزینی احتمالی مقادیر گمشده و استانداردسازی برای مدل‌های حساس به مقیاس در داخل خط لوله و فقط با داده آموزشی برازش می‌شوند. هرچند در جدول نهایی مقدار گمشده‌ای وجود ندارد، مرحله جایگزینی میانه به‌صورت دفاعی در خط لوله حفظ شده است.

برای خوشه‌بندی گفتار، هر 720 نمونه بدون استفاده از برچسب زبان پس از جایگزینی میانه، استانداردسازی و
\lr{PCA}
به 30 مؤلفه کاهش یافته‌اند:

\begin{equation}
X_{\mathrm{speech,cluster}}
\in
\mathbb{R}^{720\times30}
\end{equation}

به این ترتیب، در پایان مرحله پیش‌پردازش، هر دو مجموعه داده به فرم مناسب برای ورود به مراحل طبقه‌بندی و خوشه‌بندی تبدیل شده‌اند.


% =================================================
\section{خلاصه و جمع‌بندی}
% =================================================

در این فصل، مراحل عملی آماده‌سازی داده‌های متنی و گفتاری تشریح شد.

در بخش متن، 608 نمونه از شش زبان بارگذاری و پس از یک پاک‌سازی حداقلی، با استفاده از ویژگی‌های سطح کلمه و کاراکتر به فضای عددی تبدیل شدند. برای طبقه‌بندی، فضای ترکیبی 5000بعدی فقط روی داده آموزشی برازش و با انتخاب ویژگی به 50 بعد کاهش یافت. برای خوشه‌بندی نیز یک نمایش 50بعدی مستقل با
\lr{Truncated SVD}
و بدون استفاده از برچسب ساخته شد. تقسیم گروه‌آگاه و قرارگرفتن تمام مراحل وابسته به داده درون زنجیره آموزش، از ورود اطلاعات بخش آزمون به فرایند برازش جلوگیری می‌کند.

در بخش گفتار، 720 نمونه از چهار زبان به صورت کاملاً متوازن مورد استفاده قرار گرفتند. نگارنده فایل‌های صوتی را به قالب استاندارد تک‌کاناله با نرخ نمونه‌برداری 16 کیلوهرتز تبدیل و شدت آن‌ها را نرمال‌سازی کرده و سپس برای هر نمونه 133 ویژگی عددی استخراج کرده است. بررسی نهایی نشان داد که هیچ فایل نامعتبر، مقدار گمشده یا مقدار بی‌نهایتی در داده نهایی وجود ندارد. خط لوله آزمایش فعلی از همین جدول استخراج‌شده آغاز می‌شود و برای جداسازی آموزش و آزمون از گروه‌های گوینده یا منبع استفاده می‌کند.

خروجی مراحل این فصل، دو مجموعه داده مستقل و آماده برای ورود به الگوریتم‌های یادگیری ماشین است. نحوه آموزش مدل‌های طبقه‌بندی و خوشه‌بندی، تنظیم پارامترها و نتایج حاصل از اجرای آن‌ها در فصل بعد بررسی خواهد شد.
